\documentclass[12pt]{article}
\usepackage[T1]{fontenc}
\usepackage[utf8]{inputenc}
\usepackage{lmodern}
\usepackage{textcomp}
\usepackage{enumitem}
\usepackage{geometry}
\geometry{margin=1in}
\begin{document}
\begin{center}
Answer Key - Biology - Undergraduate Level Assessment 22 \
\small (Answers are ordered exactly as in the question paper.)
\end{center}

\section*{Multiple Choice}
\begin{enumerate}[label=\arabic*.]
\item \textbf{Answer:} D
\\ \textit{Options:} A) High metabolic rate to generate heat; B) Storage of water in fatty tissues; C) Long loops of Henle to maximize water reabsorption; D) Short loops of Henle to maximize water secretion; E) Large amounts of body hair to trap heat
\\ \textit{Note:} Short loops of Henle are associated with increased water reabsorption, which is crucial for animals in hot arid environments to conserve water, making their presence unlikely in such habitats as it would lead to excessive water loss.
\item \textbf{Answer:} A
\\ \textit{Options:} A) Courtship is primarily for entertainment purposes among terns; B) Courtship is not necessary for reproduction; C) Courtship is a way to show dominance; D) Courtship is a random behavior with no impact on reproduction; E) Courtship is necessary for physical transformation required for reproduction
\\ \textit{Note:} The correct answer highlights that courtship behaviors, such as a male tern presenting a fish, serve not only to attract a mate but also provide social bonding and entertainment, which are important for establishing pair bonds essential for successful reproduction.
\item \textbf{Answer:} A
\\ \textit{Options:} A) Oxygen is used to maintain the temperature of cells.; B) Oxygen is used by cells to produce alcohol.; C) Oxygen is used to produce carbon dioxide in cells.; D) Oxygen is solely responsible for stimulating brain activity.; E) Oxygen is used to synthesize proteins in cells.
\\ \textit{Note:} Oxygen is essential for cellular respiration, which produces energy in the form of ATP, and this energy is crucial for maintaining the metabolic processes that regulate and sustain the temperature of cells.
\item \textbf{Answer:} D
\\ \textit{Options:} A) O$_{2}$; B) CH 4 (Methane); C) CO 2; D) H$_{2}$O; E) N$_{2}$ (Nitrogen)
\\ \textit{Note:} The correct answer is H$_{2}$O because during photosynthesis, water molecules are split to release oxygen as a byproduct while the remaining hydrogen is used in the formation of glucose.
\item \textbf{Answer:} A
\\ \textit{Options:} A) Duplication of chromosome 14; B) Translocation between chromosomes 14 and 21; C) Deletion in chromosome 21; D) Mutation in chromosome 14; E) Duplication of chromosome 18
\\ \textit{Note:} The correct answer is misleading, as Down's syndrome is actually caused by the trisomy of chromosome 21, not by a duplication of chromosome 14.
\end{enumerate}

\section*{True / False}
\begin{enumerate}[label=\arabic*.]
\setcounter{enumi}{5}
\item \textbf{Answer:} False
\\ \textit{Options:} A) True; B) False
\\ \textit{Note:} The statement is false because glycolysis primarily produces NADH, ATP, and pyruvate, rather than NADP+, ADP, and sugar.
\item \textbf{Answer:} True
\\ \textit{Options:} A) True; B) False
\\ \textit{Note:} Comparing gene frequencies and genotype frequencies over two generations can indicate whether a population is in genetic equilibrium because, according to the Hardy-Weinberg principle, allele and genotype frequencies in a population that is not evolving will remain constant across generations unless influenced by evolutionary forces.
\item \textbf{Answer:} False
\\ \textit{Options:} A) True; B) False
\\ \textit{Note:} The probability of getting exactly three heads in five flips of a balanced coin is not 25 percent; it is calculated using the binomial probability formula, which yields approximately 31.25 percent.
\item \textbf{Answer:} True
\\ \textit{Options:} A) True; B) False
\\ \textit{Note:} A photosynthetic plant requires essential nutrients like iron and magnesium for chlorophyll production and overall growth; deficiencies in these nutrients lead to pale coloration, stunted growth, and weakened structure.
\item \textbf{Answer:} False
\\ \textit{Options:} A) True; B) False
\\ \textit{Note:} Cholesterol uptake primarily occurs through mechanisms such as endocytosis and transporter proteins rather than relying on adhesion plaques for cellular internalization.
\end{enumerate}

\section*{Long Form Response}
\begin{enumerate}[label=\arabic*.]
\setcounter{enumi}{10}
\item \textbf{Answer:} Hemophilia is a recessive sex-linked disorder primarily associated with the X chromosome. In a couple where their son is affected by hemophilia, it indicates that the mother is either a carrier (XhX) or affected herself (XhXh), while the father must have the normal Y chromosome (XY). Given that the son inherits his Y chromosome from his father and the affected X chromosome from his mother, the probability of their next daughter having hemophilia depends on the mother's carrier status. If the mother is a carrier, there is a 25\% chance that the daughter will inherit the affected X chromosome and thus become a carrier herself (XhX) or a 5\% chance of being affected (XhXh). If the mother is affected, the daughter will inherit the affected X chromosome and hence have a 100\% chance of being affected. Therefore, if the mother is a carrier, there is a 5\% chance that their next daughter will inherit hemophilia.
\\ \textit{Note:} correct because it accurately describes the inheritance pattern of hemophilia as a recessive sex-linked disorder, highlighting the roles of the X and Y chromosomes in determining the genetic status of the affected son and the implications for their daughter's probability of inheriting the disorder based on the mother's carrier status.
\item \textbf{Answer:} Coevolution refers to the reciprocal evolutionary changes that occur between interacting species, particularly evident in predator-prey relationships. In this dynamic, predators develop adaptations to more effectively hunt their prey, while prey species evolve defenses to evade predation, leading to a continuous cycle of adaptation and counter- adaptation. This interaction not only influences the survival and reproductive success of both groups but also has significant implications for biodiversity, as it fosters a variety of traits and behaviors within species. Additionally, coevolution impacts ecosystem dynamics by shaping community structure and influencing species interactions, ultimately contributing to the resilience and stability of ecosystems. Thus, coevolution plays a crucial role in driving evolutionary processes and maintaining ecological balance.
\\ \textit{Note:} correct because it accurately describes coevolution as a reciprocal process in predator-prey interactions, emphasizing the continuous adaptations that arise from these relationships and their broader impacts on biodiversity and ecosystem dynamics.
\end{enumerate}

\vfill
\noindent\textit{End of Answer Key}
\end{document}